\section{Utilizzo del polimorfismo}

In linea con la separazione imposta dal design pattern MVC, l'applicazione del polimorfismo
all'interno del progetto avviene in due fasi.

\subsection{Modello logico}

All'interno del modello logico, l'applicazione del polimorfismo si fonda sulla definizione di
alcune funzioni virtuali pure nella classe base astratta \texttt{attivita}.
Tali metodi forniscono un'interfaccia comune per caratteristiche trasversali ad ogni tipo di
attività (come la possibilità di essere completate), lasciando alle classi derivate l'implementazione
del comportamento specifico.
Il dominio di applicazione di queste funzioni è esclusivamente il \textbf{Model}, garantendone l'indipendenza
rispetto all'interfaccia grafica e permettendone il riutilizzo in altri contesti (cambio di framework
per l'interfaccia, passaggio da interfaccia grafica a client a linea di comando (CLI), ...).

\subsection{Design Pattern Visitor}

Per rendere la \textbf{View} il più passiva possibile rispetto ai dati in input, si è scelto di
affidare l'estrazione e la gestione dei dati grezzi al \textbf{Controller}. Alle classi
intermedie dell'interfaccia grafica viene affidato solamente il compito di inoltrare i pacchetti ricevuti
dal controller alle classi predisposte alla visualizzazione dei dati.
Per gestire i dati grezzi è stato deciso di utilizzare i DTO (Data Transfer Object), in modo da
slegare la logica delle classi dalla visualizzazione finale nell'interfaccia grafica.

L'applicazione del polimorfismo a questo processo avviene nel controller, grazie all'applicazione
del design pattern Visitor. Grazie al Double Dispatch (doppio smistamento), la classe concreta
restituisce il puntatore a se stessa \texttt{this}, invocando autonomamente la funzione corretta
del visitor.

In particolare, nel progetto i visitor sono stati usati per generare dinamicamente i DTO contenenti
i dettagli delle attività necessarie in un determinato momento all'interfaccia grafica.

Nel primo caso, ad ogni aggiornamento della schermata principale (che avviene ad ogni creazione, modifica
o eliminazione di un'attività, in modo da garantire la sincronizzazione tra \textit{Model} e
\textit{View}), il controller prende tutte le attività presenti nel \texttt{gestore}, ne genera il
DTO corrispondente e lo smista nella lista corretta, predisponendone la visualizzazione nella
schermata principale.

Invece, nel secondo caso, il visitor agisce su un'attività generica, identificata tramite ID
univoco, estraendone dinamicamente i dati grezzi e generando il rispettivo DTO, che successivamente
verrà inoltrato alla finestra richiedente (che sia essa di modifica o di visualizzazione dettagliata
in sola lettura).

Grazie all'applicazione del polimorfismo e del design pattern \textit{Visitor}, è possibile far
lavorare la classe \texttt{gestore} esclusivamente con puntatori ad attività generiche,
favorendo l'estensibilità del codice.